\section{Nothing Here Is Authorised}
\label{sec:ch15-nothing-here-is-authorised}
\index{authorization!Mosaic has none|(}%

Ask the graph for a page of products and the reviews on them, and every review
comes back with its author, and every author comes back with their email
address. No token, no header, no account. That query has worked since
chapter~\ref{ch:strangling-the-monolith} split Accounts out, and it worked
before that inside the monolith.

It is worth being precise about what is wrong there, because \enquote{the graph has
no security} is too broad to act on. The product catalogue should answer
anybody: a storefront that demands a login before it will show a price does not
sell anything. The rating and the body of a review should answer anybody too,
because that is what a review is for. One field on that path is different, and
it is the one nobody chose to publish. \code{Customer.email} exists because
Accounts needs it, and it reached the graph because it was a property on a
record that became a type.

\subsection*{Two questions, and the router keeps them apart}

\index{authentication!versus authorization}%
There are two questions and it is worth being pedantic about which is which,
because Cosmo's configuration has a block for each and they behave differently
enough that mixing them up wastes a morning.

Authentication asks who the caller is. It reads a token, checks the signature
and the claims, and produces either a principal or nothing. It refuses nobody.

Authorisation asks what that caller may have. It runs against a principal that
authentication already established, and it is the half that says no.

Keeping them apart is what lets a graph be partly public. Mosaic wants a
request with no token to succeed, and to succeed with fewer fields in the
answer than a request that has one. A single \enquote{require a login} switch
cannot express that, and the router has such a switch, which
section~\ref{sec:ch15-four-callers-one-query} turns on once to see what a graph
looks like from the other side of it.

\subsection*{Where the check goes}

\index{JWT!at the router versus at the subgraph}%
The graph is eight processes now. The router faces the client; seven services
sit behind it and face only the router, and the composer reads an eighth
subgraph that is a schema file with nothing behind it. So the token could be checked in the
router, in each subgraph, or in both, and the three arrangements are genuinely
different rather than a matter of taste.

Checking only at the router is the arrangement people describe when they draw
the diagram. One process holds the signing keys, one process talks to the
identity provider, and the subgraphs stay simple. The problem is what
\enquote{behind the router} means on a real network. Every one of Mosaic's
seven services answers on a port anybody inside the network can reach, and
chapter~\ref{ch:the-first-cut-extracting-catalog} spent a section testing a
subgraph directly with Postman because that is a normal thing to do. A rule
that lives only in the router is a rule that a single \code{curl} skips.

Checking only in the subgraphs puts the rule next to the data, which is where
it can be most precise: Accounts is the only process that knows what
\code{Customer.email} means. The cost is that the router now plans a query
without knowing which parts of it will be refused, so it fetches and then finds
out, and every service needs the signing key.

Mosaic ends this chapter doing both, which is the honest answer and is also the
one that exposes the thing this chapter is really about. The two mechanisms
HotChocolate gives you for these two positions do not know about each other,
and the seam between them is where the next section starts.

I should say plainly what this chapter does not do. It puts one scope on one
field and one login requirement on one root field. It does not write a policy,
it does not use roles, and it does not authorise \code{Query.node}, which
chapter~\ref{ch:hard-modeling-problems} handed here as a named risk and which
section~\ref{sec:ch15-what-the-composer-checks} explains rather than solves.
What I wanted from Mosaic here was the smallest change that makes the seam
visible, because the seam is the part that is not in anybody's documentation.

\medskip

That seam is two attributes, and they look interchangeable until you export the
schema.

\index{authorization!Mosaic has none|)}%
